% Section 10.7, translated from sv/chapters/10/exercises.tex.

\section{Exercises}
\label{c10:sec:uppgifter}

\begin{aside}
\textbf{How to work with the exercises.} Parts~1 and~2 are done during the lesson: first on your
own, a few minutes per exercise, then together. Part~1 reviews the trigonometry of \kapref{9}, and
Part~2 is the chapter's new material. Part~3 is extra exercises for further practice after the
lesson.

\facitnotis{L10}
\end{aside}

% ------------------------------------------------------------------------------------------
\uppgiftsdel{Part 1}{Review exercises}

\uppgift{1.1}{Angles to radians}
Convert the following angles to radians.
\begin{delar}(4)
  \task $15^{\circ}$
  \task $-75^{\circ}$
\end{delar}

\uppgift{1.2}{Angles to degrees}
Convert the following angles to degrees.
\begin{delar}(4)
  \task $-\pi/6$
  \task $3\pi/2$
\end{delar}

\uppgift{1.3}{Determining the properties of an AC voltage}
An AC voltage is shown in the figure below.

\bild[0.75\linewidth]{lectures/L10/appendix/images/1.3_sine_wave.png}

Find the voltage's equation in the form $u(t) = \abs{U}\sin(\omega t + \delta)$.

\uppgift{1.4}{Equation and graph of an AC voltage}
An AC voltage has the amplitude $4\,\text{V}$, the frequency $50\,\text{Hz}$ and the phase
$-30^{\circ}$. Find the equation $u(t)$ of the AC voltage, with the phase in radians, and draw the
sine curve over one period $T$.

\uppgift{1.5}{Calculating the phase of an alternating current}
The equation of a given AC voltage is
\[
  u(t) = 6\sin(80\pi t + \delta).
\]
At the time $t = 15\,\text{ms}$, $u(t) = 3\,\text{V}$. Calculate the phase $\delta$.

% ------------------------------------------------------------------------------------------
\uppgiftsdel{Part 2}{New material}

\uppgift{2.1}{Logarithmic equations}
Solve the following logarithmic equations.
\begin{delar}(3)
  \task $5^x = 125$
  \task $2^{x-3} = 128$
  \task $e^{x+1} = 120$
\end{delar}

\uppgift{2.2}{Half-life of the charge in a battery}
The charge in a battery can be described by the corresponding voltage with the function
\[
  u(t) = U_0 \cdot a^t,
\]
where
\begin{itemize}
  \item $u(t)$ is the present charge,
  \item $U_0$ is the original charge,
  \item $a$ is the growth factor,
  \item $t$ is the number of hours that have passed.
\end{itemize}
A battery loses half its charge in $40$ hours. Calculate after how long only $10\,\%$ of the
charge remains.

\uppgift{2.3}{Linear gain}
Two signals have the voltage levels $L_1 = 15\,\text{dB}$ and $L_2 = 49\,\text{dB}$. Calculate the
linear voltage gain $G_{\text{lin}}$ between these two levels. Use the formula
\[
  G_{\text{dB}} = 20\log_{10}(G_{\text{lin}})
\]
and remember that $G_{\text{dB}} = L_2 - L_1$.

\uppgift{2.4}{Calculating an RMS value in dBV}
The relationship between the RMS value of a voltage and the corresponding value in dBV (decibels
relative to one volt) is
\[
  U_{\text{dBV}} = 20\log_{10}\left(\frac{U_{\text{RMS}}}{1\,\text{V}}\right),
\]
where $U_{\text{RMS}}$ is the RMS value of the voltage in V and $U_{\text{dBV}}$ the voltage in
dBV. A sinusoidal voltage has the level $17.0\,\text{dBV}$. Find the amplitude in V.

% ------------------------------------------------------------------------------------------
\uppgiftsdel{Part 3}{Extra exercises}
The exercises below are extra practice, best done on your own after the lesson. Most of them can
be done in your head or with pencil and paper.

\uppgift{3.1}{Logarithms without a calculator}
Calculate.
\begin{delar}(6)
  \task $\log 100$
  \task $\log 1000$
  \task $\log 1$
  \task $\log 0.01$
  \task $\ln e$
  \task $\ln 1$
\end{delar}

\uppgift{3.2}{The laws of logarithms}
Write as a single logarithm and simplify as far as possible.
\begin{delar}(3)
  \task $\log 3 + \log 4$
  \task $\log 20 - \log 4$
  \task $2\log 5$
  \task $\log 2 + \log 5$
  \task $3\log 2 - \log 4$
\end{delar}

\uppgift{3.3}{Exponential equations}
Solve the equations.
\begin{delar}(5)
  \task $2^x = 32$
  \task $10^x = 0.001$
  \task $5^x = 100$
  \task $e^x = 20$
  \task $4^{x+1} = 64$
\end{delar}

\uppgift{3.4}{Logarithmic equations}
Solve the equations.
\begin{delar}(3)
  \task $\log x = 2$
  \task $\ln x = 0$
  \task $\log(x + 5) = 1$
  \task $2\log x = 4$
  \task $\ln(2x) = 1$
\end{delar}

\uppgift{3.5}{Voltage gain in dB}
Use
\[
  G_{\text{dB}} = 20\log_{10}\left(\frac{U_{\text{out}}}{U_{\text{in}}}\right).
\]
\begin{parts}
  \item $U_{\text{in}} = 10\,\text{mV}$ and $U_{\text{out}} = 1\,\text{V}$. Calculate
    $G_{\text{dB}}$.
  \item $U_{\text{in}} = 2\,\text{V}$ and $U_{\text{out}} = 2\,\text{V}$. Calculate
    $G_{\text{dB}}$.
  \item $U_{\text{in}} = 5\,\text{V}$ and $U_{\text{out}} = 0.5\,\text{V}$. Calculate
    $G_{\text{dB}}$.
  \item What does a negative value of $G_{\text{dB}}$ mean?
\end{parts}

\uppgift{3.6}{From dB to linear gain}
Use
\[
  G_{\text{lin}} = 10^{G_{\text{dB}}/20}.
\]
\begin{delar}(4)
  \task $G_{\text{dB}} = 20\,\text{dB}$
  \task $G_{\text{dB}} = 60\,\text{dB}$
  \task $G_{\text{dB}} = 6\,\text{dB}$
  \task $G_{\text{dB}} = -3\,\text{dB}$
  \task* An amplifier with $G_{\text{dB}} = 26\,\text{dB}$ is fed with
    $U_{\text{in}} = 50\,\text{mV}$. Calculate $U_{\text{out}}$.
\end{delar}

\uppgift{3.7}{Power gain in dB}
Use
\[
  G_{\text{dB}} = 10\log_{10}\left(\frac{P_{\text{out}}}{P_{\text{in}}}\right).
\]
\begin{parts}
  \item $P_{\text{in}} = 1\,\text{W}$ and $P_{\text{out}} = 100\,\text{W}$. Calculate
    $G_{\text{dB}}$.
  \item $P_{\text{in}} = 20\,\text{W}$ and $P_{\text{out}} = 10\,\text{W}$. Calculate
    $G_{\text{dB}}$.
  \item Why is the factor $10$ used for power but $20$ for voltage?
  \item How many dB does a doubling of the power correspond to? And a doubling of the voltage?
\end{parts}

\uppgift{3.8}{Cascaded stages}
A signal chain consists of two amplifier stages with $G_1 = 12\,\text{dB}$ and
$G_2 = 18\,\text{dB}$, and a cable that attenuates $4\,\text{dB}$.
\begin{parts}
  \item Calculate the total gain in dB.
  \item Calculate the total linear gain.
  \item Calculate each stage linearly and check that the product agrees with b).
  \item Why are decibels convenient to calculate with when stages are cascaded?
\end{parts}

\uppgift{3.9}{Level in dBV}
Use
\[
  U_{\text{dBV}} = 20\log_{10}\left(\frac{U_{\text{RMS}}}{1\,\text{V}}\right).
\]
\begin{parts}
  \item $U_{\text{RMS}} = 1\,\text{V}$. Calculate the level in dBV.
  \item $U_{\text{RMS}} = 10\,\text{V}$. Calculate the level in dBV.
  \item $U_{\text{RMS}} = 0.1\,\text{V}$. Calculate the level in dBV.
  \item A signal is at $12\,\text{dBV}$. Calculate $U_{\text{RMS}}$.
  \item Calculate the amplitude $\abs{U}$ of the signal in d).
\end{parts}

\uppgift{3.10}{Half-life and doubling time}
A quantity is described by $f(t) = C \cdot a^t$.
\begin{parts}
  \item A battery loses $2\,\%$ of its charge per hour. Find the growth factor $a$.
  \item After how many hours is half the charge left?
  \item A set of measurement data grows by $10\,\%$ per day. After how many days has it doubled?
  \item Show in general that the doubling time is given by $t = \dfrac{\log 2}{\log a}$.
\end{parts}

\uppgift{3.11}{Find the error}
Each line contains a common mistake. Explain the mistake and give the correct answer.
\begin{parts}
  \item $\log(a + b) = \log a + \log b$
  \item $\dfrac{\log 8}{\log 2} = \log 4$
  \item $\log(a^n) = (\log a)^n$
  \item $G_{\text{dB}} = 0$ means that the signal disappears completely.
  \item $10^{\log x} = 10x$
  \item $\ln 0 = 1$
\end{parts}
